% policy-document.tex — LaTeX equivalent of templates/typst/policy-document.typ.
%
% Unlike the Typst template, this one does NOT read the policy JSON: LaTeX has
% no native data loading, and the alternatives (generating .tex from a script,
% or driving LuaTeX) are the error-prone path DOCUMENT_TEMPLATES.md argues
% against. Fill the fields in by hand, or generate this file from the JSON with
% a small script if you need it automated.
%
% Build:
%   cd templates/latex && latexmk -xelatex policy-document.tex
%   (or: xelatex policy-document.tex, twice, for the page-count reference)

\documentclass[11pt,a4paper]{article}
\usepackage{grc}

% ---- Document metadata ----------------------------------------------------
\setcldoctitle{Access Control Policy}
\setclreference{POL-AC-001}
\setclclassification{Confidential}

\begin{document}

\clcoverpage{Access Control Policy}{Policy}{Northwind Financial Services Ltd}

% Ragged right with hyphenation suppressed: these are short field values in
% narrow columns, where a hyphenated break ("Of-ficer") is far more visible than
% a ragged edge.
\begin{tabularx}{\linewidth}{@{}>{\raggedright\hyphenpenalty=10000\arraybackslash}X>{\raggedright\hyphenpenalty=10000\arraybackslash}X>{\raggedright\hyphenpenalty=10000\arraybackslash}X@{}}
  \cllabel{Reference}\newline POL-AC-001 &
  \cllabel{Document type}\newline Policy &
  \cllabel{Version}\newline v2.0 \\[11pt]
  \cllabel{Document owner}\newline Head of Information Security &
  \cllabel{Approver}\newline Chief Information Security Officer &
  \cllabel{Effective date}\newline 2026-02-01 \\[11pt]
  \cllabel{Review cycle}\newline 12 months &
  \cllabel{Next review}\newline 2027-02-01 &
  \cllabel{Frameworks}\newline ISO 27001, NIST 800-53 \\
\end{tabularx}

\clearpage
\pagestyle{clbody}

% ---- Document control -----------------------------------------------------
\section*{Document Control}
\addcontentsline{toc}{section}{Document Control}

\begin{clkvtable}
  Reference & POL-AC-001 \\
  Title & Access Control Policy \\
  Document type & Policy \\
  Classification & Confidential \\
  Status & Approved \\
  Document owner & Head of Information Security \\
  Approver & Chief Information Security Officer \\
  Client & Northwind Financial Services Ltd \\
  Frameworks & ISO 27001, NIST 800-53, PCI DSS \\
  Effective date & 2026-02-01 \\
  Next review & 2027-02-01 \\
\end{clkvtable}

\vspace{1em}
\begin{clcallout}
  \cllabel{Summary}\par\vspace{0.2em}
  This policy establishes the requirements for granting, reviewing and revoking
  access to Northwind information systems and the cardholder data environment.
  It applies to all employees, contractors and third parties, and is mandatory
  for any system that processes, stores or transmits regulated data.
\end{clcallout}

% ---- Revision history -----------------------------------------------------
% A clause 7.5 requirement, not a nicety: an auditor asks what changed between
% versions and who approved it.
\vspace{1.2em}
\section*{Revision History}
\addcontentsline{toc}{section}{Revision History}

\rowcolors{2}{}{clTableZebra}
\begin{tabularx}{\linewidth}{@{}>{\sffamily\small}l>{\sffamily\small}l>{\small}lX@{}}
  \toprule
  \sffamily\bfseries\small Version & \sffamily\bfseries\small Approved &
  \sffamily\bfseries\small Approved by & \sffamily\bfseries\small Change summary \\
  \midrule
  v2.0 & 2026-02-01 & Chief Information Security Officer &
    Added quarterly privileged access review and one-business-day revocation requirement. \\
  v1.0 & 2025-01-15 & Chief Information Security Officer & Initial issue. \\
  \bottomrule
\end{tabularx}

\clearpage
\renewcommand{\contentsname}{Contents}
\tableofcontents
\clearpage

% ---- Body -----------------------------------------------------------------
\section{Purpose}
Unauthorised access is the most frequently exploited weakness in financial
services environments, and the majority of material incidents begin with a
credential that should have been revoked or was never appropriate to issue.
This policy sets the mandatory requirements that limit access to information
systems to identified, authorised individuals, and that ensure access is removed
promptly when it is no longer justified.

It exists to protect the confidentiality, integrity and availability of
Northwind information assets, and to meet the access control obligations placed
on the organisation by ISO/IEC 27001:2022, NIST SP 800-53 Revision 5 and the
Payment Card Industry Data Security Standard.

\section{Scope}
This policy applies to all production and pre-production information systems
operated by or on behalf of Northwind, including the cardholder data
environment, and to every account type used to access them --- interactive user
accounts, privileged administrative accounts, service accounts and accounts held
by third parties.

\section{Policy Statements}
% Normative language only: must / shall (mandatory), should (expected, deviation
% needs justification), may (discretionary). An assessor cannot test "will",
% "strives to" or "is encouraged to".
Access \textbf{must} be granted on the principle of least privilege: an account
must hold only the permissions required for the role it supports, and no more.

Access \textbf{must} be authorised by the system owner before it is provisioned,
and the authorisation must be recorded and retained for a minimum of twelve
months.

Privileged access \textbf{must} be issued separately from an individual's
standard account, \textbf{must} require multi-factor authentication, and
\textbf{must} be reviewed at least quarterly by the system owner.

Access \textbf{must} be revoked within one business day of a person leaving the
organisation or changing to a role that no longer requires it.

\section{Roles and Responsibilities}
The Chief Information Security Officer owns this policy, approves it, and is
accountable for its effectiveness. The Head of Information Security must
maintain this policy and report material non-compliance to the Executive Risk
Committee. System owners must authorise access to the systems they own and
conduct the periodic access reviews required by this policy.

\section{Management Commitment}
The Executive Risk Committee endorses this policy and commits the resources
required to implement and sustain it. Compliance with this policy is a condition
of continued access to Northwind systems.

\section{Compliance and Enforcement}
Compliance is monitored through the periodic access reviews, automated reporting
on privileged account usage, and internal audit. Breaches must be reported to
the Head of Information Security and are handled under the disciplinary
procedure.

% ---- Control mapping ------------------------------------------------------
\clearpage
\section{Control Mapping}
This document is mapped to the following catalog controls. Coverage is stated
per control: \emph{full} means this document satisfies the control on its own,
\emph{partial} means further documents are required, and \emph{supporting}
means the text contributes context without satisfying the control.

\vspace{0.6em}
\rowcolors{2}{}{clTableZebra}
\begin{tabularx}{\linewidth}{@{}>{\ttfamily\small}lX>{\small}l>{\small}l@{}}
  \toprule
  \sffamily\bfseries\small Control & \sffamily\bfseries\small Name &
  \sffamily\bfseries\small Coverage & \sffamily\bfseries\small Section \\
  \midrule
  AC-1 & Policy and Procedures & Full & Purpose \\
  AC-2 & Account Management & Full & Policy Statements \\
  AC-3 & Access Enforcement & Full & Policy Statements \\
  AC-6 & Least Privilege & Partial & Policy Statements \\
  \bottomrule
\end{tabularx}

\vspace{1.5em}
\clapproval{Chief Information Security Officer}{Head of Information Security}{2026-02-01}

\end{document}
